% =====================================================================
%  2bat-ud3-15.tex  —  SESSIÓ 15: Taller de preparació de la prova
% =====================================================================
\horatitol{Sessió 15 — Taller de preparació de la prova}%
{Conèixer l'estructura de la prova i la rúbrica, assajar un simulacre per blocs i fer un pla d'estudi personal.}

\subsection{Idea clau}

Ens preparem per a la prova \textbf{sense sorpreses}. La prova de la sessió següent té
\textbf{quatre blocs}, un per cada criteri (C1-C4). Saber exactament què es demana i
com es corregeix fa que l'examen deixi de ser imprevisible. Després assajarem amb un
\textbf{simulacre}.

\begin{idea}
\textbf{Estructura de la prova} (un bloc per criteri):
\begin{enumerate}[leftmargin=1.6em, itemsep=2pt]
  \item \textbf{Bloc 1 (C1) — Límits a l'infinit:} estimar i interpretar el límit d'una
        funció (des d'una taula, gràfica o expressió).
  \item \textbf{Bloc 2 (C2) — Asímptotes:} identificar les asímptotes d'una funció i
        dir-ne el significat.
  \item \textbf{Bloc 3 (C3) — Continuïtat:} comprovar si una funció a trossos és
        contínua o fa un salt a una frontera.
  \item \textbf{Bloc 4 (C4) — Lectura global:} descriure i interpretar el comportament
        complet d'un model social.
\end{enumerate}
\textbf{Es valora}, a més del resultat: la \textbf{interpretació} en context, la
\textbf{comprovació acurada} de la continuïtat i la \textbf{lectura social} del model.
\end{idea}

\subsection{Rúbrica simplificada}

\noindent Així es corregeix cada bloc (nivells d'assoliment):
\begin{center}
\renewcommand{\arraystretch}{1.3}
\begin{tabular}{p{2.4cm} p{10.2cm}}
\toprule
\textbf{Nivell} & \textbf{Què demostra} \\
\midrule
\textbf{AE} (excel·lent) & Resol sense errors i interpreta el resultat en el context
social. \\
\textbf{AN} (notable) & Resol correctament amb algun error menor; la interpretació és
essencialment bona. \\
\textbf{AS} (suficient) & Resol el procediment bàsic amb alguns errors; interpretació
incompleta. \\
\textbf{NA} (no assolit) & No identifica el procediment o comet errors que
n'invaliden el resultat. \\
\bottomrule
\end{tabular}
\end{center}

\subsection{Simulacre — un ítem de cada bloc}

\begin{exemple}[com es resol un ítem ben fet]
\textbf{Ítem tipus (Bloc 3):} comprova si
$f(x)=\begin{cases}2x+1 & x\le 2\\ x+3 & x>2\end{cases}$ és contínua a $x=2$. Un
\textbf{ítem complet} avalua els dos trams i conclou:
\[
\text{esquerre } 2\cdot 2+1=5,\quad \text{dret } 2+3=5 \ \Rightarrow\ \text{contínua}.
\]
\emph{Avaluar tots dos costats i concloure clarament distingeix un AE d'un AS.}
\end{exemple}

\begin{exercicis}
\textbf{Resol el simulacre amb temps controlat. Un ítem per bloc:}
\begin{exlist}
  \item \textbf{(Bloc 1 — C1)} Estima $\displaystyle\lim_{x\to\infty}\left(400-\frac{800}{x}\right)$
        amb una taula ($x=10,100,1000$) i interpreta el resultat si és l'ocupació d'un
        servei.
  \item \textbf{(Bloc 2 — C2)} Per a $f(x)=\dfrac{30}{x}$, digues les dues asímptotes
        (horitzontal i vertical) i interpreta-les si $x$ és el nombre de beneficiaris
        d'un fons.
  \item \textbf{(Bloc 3 — C3)} Comprova si
        $g(x)=\begin{cases}x^2 & x\le 2\\ 3x-2 & x>2\end{cases}$ és contínua a $x=2$ i,
        si no, digues la mida del salt.
  \item \textbf{(Bloc 4 — C4)} Un model creix els primers mesos, s'acosta a una
        asímptota $y=500$ i al mes $10$ fa un salt cap avall fins a $350$. Descriu-ne el
        comportament global i què implica socialment.
\end{exlist}
\end{exercicis}

\subsection{Formulari-resum (suport)}

\begin{idea}
\textbf{Tingues-ho a mà mentre prepares la prova:}
\begin{itemize}[leftmargin=1.4em, itemsep=2pt]
  \item \textbf{Límit a l'infinit:} valor cap al qual s'acosta $f(x)$ quan $x$ creix
        molt (taula o gràfica); si $\frac{k}{x}$ apareix, tendeix a $0$.
  \item \textbf{Asímptota horitzontal $y=L$:} correspon a $\lim_{x\to\infty}f(x)=L$
        (sostre de saturació).
  \item \textbf{Asímptota vertical $x=a$:} valor que la funció no pot prendre (p.\,ex.\
        denominador zero), on es dispara.
  \item \textbf{Continuïtat a $x=a$:} avalua els dos trams a $x=a$; si coincideixen,
        contínua; si no, salt (mida = diferència).
\end{itemize}
\end{idea}

\noindent\textbf{Pla d'estudi personal.} A partir de la teva graella d'autoavaluació
(sessió 7) i del simulacre, anota els \textbf{dos o tres aspectes} que has de repassar
amb més atenció:
\begin{center}
\renewcommand{\arraystretch}{1.5}
\begin{tabular}{c p{10.5cm}}
\toprule
\textbf{\#} & \textbf{Què he de repassar (i com)} \\
\midrule
1 & \rule{9.8cm}{0.4pt} \\
2 & \rule{9.8cm}{0.4pt} \\
3 & \rule{9.8cm}{0.4pt} \\
\bottomrule
\end{tabular}
\end{center}

% ---------------------------------------------------------------------
\fullsolucions{Sessió 15 — simulacre}
\begin{solucions}
\begin{sollist}
  \item \textbf{(Bloc 1)} $f(10)=320$, $f(100)=392$, $f(1000)=399,2$: s'acosta a $400$.
        $\lim_{x\to\infty}f(x)=400$: el servei s'estabilitza en $400$ (la seva
        capacitat).
  \item \textbf{(Bloc 2)} Asímptota horitzontal $y=0$ (amb molts beneficiaris l'ajut
        tendeix a $0$, es dilueix) i vertical $x=0$ (amb gairebé cap beneficiari l'ajut
        es dispara, situació impossible).
  \item \textbf{(Bloc 3)} A $x=2$: esquerre $2^2=4$; dret $3\cdot 2-2=4$. Coincideixen:
        \textbf{contínua}.
  \item \textbf{(Bloc 4)} El servei creix i s'estabilitza acostant-se a $500$
        (asímptota: la seva capacitat); al mes $10$ un salt cap avall fins a $350$
        (discontinuïtat: una retallada o canvi brusc) que redueix de cop el servei en
        $150$ places, i a partir d'aleshores es manté en $350$. Socialment, unes $150$
        persones perden l'atenció de sobte.
\end{sollist}
\end{solucions}
